% !TEX TS-program = lualatex
% !TEX encoding = UTF-8 Unicode
\documentclass[ngerman,10pt]{article}
\usepackage{fontspec}
\usepackage[ngerman]{babel}
\usepackage[a4paper, total={16cm, 26cm}]{geometry}
\usepackage{amsmath}
\usepackage{amsfonts}
\usepackage{amssymb}
\usepackage{multicol}
\usepackage{tcolorbox}
\usepackage[ngerman]{datetime2}
\usepackage{enumitem}
\usepackage{url}
\usepackage[colorlinks=true, linkcolor=blue, urlcolor=blue]{hyperref}
\usepackage{listings}

\lstset{
    language=Python,
    basicstyle=\ttfamily\small,
    keywordstyle=\color{blue},
    stringstyle=\color{red},
    commentstyle=\color{gray},
    breaklines=true,
    frame=single,
    framesep=1mm,
    framextopmargin=1mm,
    framexbottommargin=1mm,
    showspaces=false,
    showstringspaces=false,
    numbers=left,
    numberstyle=\tiny\color{gray},
    literate=
      {ä}{{\"a}}1 {ö}{{\"o}}1 {ü}{{\"u}}1
      {Ä}{{\"A}}1 {Ö}{{\"O}}1 {Ü}{{\"U}}1
      {ß}{{\ss}}1
}
\usepackage{titlesec}
\titleformat{\section}{\large\bfseries}{}{0em}{}

\date{\DTMgermanmonthname{\month} \the\year}

\begin{document}
\shorthandoff{"}
\thispagestyle{empty}
\hrule
\begin{center}
{\large\bf Python: Schleifen - Ausgewählte Aufgaben}\\[-0mm]
Studienkolleg der TU Berlin\hfill Till Zoppke\\
Informatik \hfill \DTMgermanmonthname{\month} \the\year\\[2mm]
\end{center}
\hrule
%\par\bigskip

% Aufgabe 1
\section*{Aufgabe 1: Eingabevalidierung}
Schreiben Sie ein Python Programm. Das Programm soll die Benutzerin so lange nach einer Zahl fragen, bis er eine gültige Zahl zwischen 1 und 10 eingibt. Für ungültige Eingaben (zu klein, zu groß, keine Zahl) soll der Benutzerin eine entsprechende Nachricht angezeigt werden. Abschließend soll das Quadrat der Zahl ausgegeben werden.

% Aufgabe 2
\section*{Aufgabe 2: break und continue}
Führen Sie das folgende Programm gedanklich durch. Geben Sie an, was das Programm ausgibt.
\begin{lstlisting}
zahlen = [4, 7, 2, 9, 1, 6, 3]
summe = 0
for z in zahlen:
    if z % 2 == 0:
        print(f"Überspringe {z}")
        continue
    if summe + z > 15:
        print(f"Abbruch bei {z}")
        break
    summe = summe + z
    print(f"Addiere {z}, Summe = {summe}")
print(f"Ergebnis: {summe}")
\end{lstlisting}

% Aufgabe 3
\section*{Aufgabe 3: Schritt-für-Schritt Ausführung}
\begin{enumerate}[label=\alph*)]

\item Führen Sie das folgende Programm gedanklich durch und geben Sie für jeden Schleifendurchlauf an, was in Zeile 3 ausgegeben wird.
\begin{lstlisting}
for i in range(2):
    for j in range(3):
        print(f"{i} x {j} = {i * j}")
\end{lstlisting}

\item Führen Sie das folgende Programm gedanklich durch. Geben Sie für jeden Schleifendurchlauf an, was in den Zeilen 4 und 6 ausgegeben wird.
\begin{lstlisting}
liste = [8, 3, 12, 1]
while liste:
    e = liste.pop(0)
    print(f"Groß: {e}") if e > 5 else print(f"Klein: {e}")
    if e > 5:
        liste.append(e // 2)
    print(f"Liste: {liste}")
\end{lstlisting}

\end{enumerate}

% Aufgabe 4
\section*{Aufgabe 4: Mittelwert berechnen}
Schreiben Sie eine Funktion \texttt{mittelwert(liste)}, die den Mittelwert einer Liste von Zahlen berechnet und zurückgibt. Verwenden Sie eine \texttt{for}-Schleife, um die Summe aller Elemente zu berechnen. Verwenden Sie \emph{nicht} die eingebauten Funktionen \texttt{sum()} oder \texttt{len()}.

\begin{lstlisting}
def mittelwert(liste):
    # Ihr Code hier

messwerte = [3.5, 7.2, 1.8, 9.4, 5.1]
ergebnis = mittelwert(messwerte)
print(f"Mittelwert: {ergebnis}") # Erwartete Ausgabe: Mittelwert: 5.4
\end{lstlisting}

\end{document}
